\documentclass[12pt]{article}

\usepackage{fullpage}
\usepackage{multicol,multirow}
\usepackage{tabularx}
\usepackage{ulem}
\usepackage[utf8]{inputenc}
\usepackage{mathtools}
\usepackage{pgfplots}
\usepackage[russian]{babel}

\begin{document}

\section*{Лабораторная работа №\,7 по курсу дискрeтного анализа: Динамическое программирование}

Выполнил студент группы М8О-308Б-22 МАИ \textit{Немкова Анастасия}.

\subsection*{Условие}
Задано целое число n. Необходимо найти количество натуральных (без нуля) чисел, которые меньше n по значению и меньше n лексикографически (если сравнивать два числа как строки), а так же делятся на m без остатка.

\textbf{Формат ввода}

В первой строке строке задано 1 <= n <= $10^{18}$ и  1 <= m <= $10^5$.

\textbf{Формат вывода}

Необходимо вывести количество искомых чисел.


\subsection*{Метод решения}

Динамическое программирование — это подход, основанный на разбиении сложной задачи на более простые подзадачи. 

В этой задаче мы должны посчитать количество чисел, которые не превышают заданное число n и делятся на m. Основная идея заключается в том, что мы будем поочередно обрабатывать каждую цифру числа n и определять, сколько чисел можно сформировать, соблюдая ограничения по разрядам.

Мы используем трехмерный массив dp, где:
\begin{itemize}
    \item Первый индекс представляет текущую позицию цифры в числе n.
    \item Второй индекс фиксирует остаток от деления на m.
    \item Третий индекс указывает, меньше ли сформированное число, чем текущее префиксное число из n (для того, чтобы контролировать, можем ли мы использовать любую цифру от 0 до 9 или ограничены текущей цифрой).
\end{itemize}

Изначально мы устанавливаем значение dp[0][0][0] = 1, что означает, что существует один способ создать пустое число с нулевым остатком. Затем мы проходим по каждой позиции в числе n и рассматриваем возможные значения для текущего разряда. Если мы не превысили число n, мы можем использовать любую цифру от 0 до 9, в противном случае нам нужно ограничиться текущей цифрой в n.

В конце мы суммируем все значения в массиве dp, где остаток равен нулю, чтобы получить количество чисел, которые соответствуют условиям задачи.

\subsection*{Описание программы}

В программе реализована функция countNumber, которая подсчитывает количество натуральных чисел, меньше заданного числа n по значению и лексикографически, а также делящихся на m без остатка. Для этого используется трехмерный вектор dp размером ((length(n) + 1) * m * 2), где:

length(n) — длина числа n,

m — модуль для деления,

2 — флаг, указывающий, меньше ли текущее число, чем соответствующий префикс числа n.

Функция инициализирует начальное состояние, затем с помощью трех вложенных циклов заполняет dp, перебирая возможные позиции, остатки и состояния. Временная сложность алгоритма составляет O(length(n) * m * 20)

\subsection*{Дневник отладки}

\begin{enumerate}
    \item 23 окт 2024, 18:44:45 WA на 21 тесте
		
    Переполнение типа. Решение: изменить тип данных результата на long long 


\end{enumerate}

\subsection*{Тест производительности}

Для измерения производительсти сравним время работы данного алгоритма и наивного, в котором перебираем все числа от m до n, увеличивая каждое на m и  для каждого проверяем, меньше ли оно n лексикографически. Протестируем на входных данных, где n = $10^k$ - 1, m = 1, k = 1, ..., 5.
	
\begin{figure}[htbp]
    \centering
    \begin{tikzpicture}
        \begin{axis}[
            xlabel={Объём входных данных ($\times 10^2$)},
            ylabel={Время работы (ms)},
            grid=major,
            xmin=0, xmax=1000,
            ymin=0, ymax=10,
            xtick={0, 200, 400, 600, 800, 1000},
            ytick={0, 2, 4, 6, 8, 10},
            legend style={at={(0.5,-0.2)},anchor=north},
            legend columns=2,
            width=0.8\textwidth,
            height=0.5\textwidth,
            ]
            \addplot[color=blue,mark=*] coordinates {
                (0.09, 0.066)
                (0.9, 0.057)
                (9, 0.056)
                (99, 0.079)
                (999, 0.08)

            };
            \addlegendentry{Мой алгоритм}
            \addplot[color=red,mark=square] coordinates {
                (0.09, 0.035)
                (0.9, 0.056)
                (9,  0.167)
                (99, 1.21)
                (999, 9.28)
            };
            \addlegendentry{Наивный}
        \end{axis}
    \end{tikzpicture}
    \label{fig:graph}
\end{figure}

\newpage
\subsection*{Выводы}

В ходе данной лабораторной работы был реализован алгоритм динамического программирования для подсчета чисел, удовлетворяющих заданным условиям, что позволило сократить вычислительные затраты по сравнению с наивным методом.

Эффективность данного подхода проявляется в возможности решения задач с множеством вариантов выбора, таких как комбинаторные задачи, задачи о рюкзаке и оптимизации. В результате, динамическое программирование не только улучшает производительность программ, но и предоставляет более структурированный способ решения сложных проблем, делая их более понятными и управляемыми.

\end{document}




